
Experiments were designed to answer three research questions corresponding to framework mechanisms: (RQ1) Does micro-pilot overhead correlate with full-scale overhead? (RQ2) Do Bayesian updates reduce prediction error? (RQ3) Does Gate 1 achieve >80\% viability classification accuracy? Each question tests a specific mechanism (M1, M2, M3) with quantified success criteria and statistical validation.

\subsubsection{4.1 Research Questions}

\textbf{RQ1 (Mechanism M1):} Does micro-pilot overhead $O_{10}$ exhibit strong correlation (r > 0.7) with full-scale overhead O\_full across diverse hypothesis types?

Linear extrapolation O\_full = k $\times O_{10}$ requires predictive scaling. If r < 0.7, correlation is too weak for reliable extrapolation.

\textbf{RQ2 (Mechanism M2):} Do Bayesian updates combining Gate 1 prior with Gate 2 likelihood reduce prediction error by >40\% compared to Gate 1 alone?

Bayesian refinement adds complexity (100-sample investment at Gate 2). If error reduction is marginal (<20\%), the framework should default to Gate 1-only prediction.

\textbf{RQ3 (Mechanism M3):} Does Gate 1 viability classification achieve >80\% accuracy for predicting whether hypotheses exceed an overhead threshold, compared to 50\% random baseline?

This validates the primary claim. Accuracy $\leq 60$\% would indicate performance no better than random guessing plus small margin.

\subsubsection{4.2 Corpus Construction}

A synthetic corpus of 32 ML hypotheses was constructed, each with both micro-pilot (10-sample) and full-scale overhead measurements. Each hypothesis represents a variation in model architecture (attention mechanisms, gradient penalties, regularization techniques, normalization schemes) evaluated on standard benchmarks.

\textbf{Rationale for Synthetic Data:} Retrospective validation requires published papers reporting micro-pilot timing data, a constraint rarely satisfied in ML literature. Papers with Code and conference proceedings (NeurIPS, ICML, ICLR) typically report full-scale results but omit 10-sample ablations. Rather than forgo validation, synthetic data was generated preserving statistical properties to be tested (scaling correlation, Bayesian error reduction, classification accuracy) while controlling for confounds. This establishes proof-of-concept for framework mechanics. External validity testing on real published papers is identified as critical future work.

\textbf{Corpus Structure:} 32 hypotheses stratified across overhead levels:
\begin{itemize}
\item 11 low-overhead (<20\% full-scale overhead)
\item 11 mid-overhead (20-80\%)
\item 10 high-overhead (>80\%)
\end{itemize}

Stratification ensures balanced representation. Coefficient of variation CV = 0.044 < 0.5 threshold confirms adequate distribution balance.

\textbf{Overhead Threshold:} T = 10\% reflects real-time deployment constraints where overhead must remain minimal. This strict threshold results in 26 of 32 hypotheses (81.25\%) classified as non-viable, a realistic distribution for deployment-constrained scenarios.

\textbf{Data Generation:} For each hypothesis:
\begin{enumerate}
\item Base overhead sampled from stratified distribution
\item Micro-pilot overhead $O_{10}$ generated with scaling factor k=1.000 (perfect linearity for proof-of-concept)
\item Mid-scale overhead $O_{100}$ = O\_full / (dataset\_size / 100)
\item Full-scale overhead O\_full computed deterministically from base overhead
\end{enumerate}

This construction enforces perfect linear scaling (r=1.000) by design, representing idealized conditions. Real-world validation will test whether r $\geq$ 0.7 threshold holds under measurement noise and hardware variance.

\subsubsection{4.3 Evaluation Metrics}

\textbf{RQ1 Metrics (Correlation):}
\begin{itemize}
\item \textbf{Pearson r:} Linear correlation between $O_{10}$ and O\_full. Success: r > 0.7, p < 0.05.
\item \textbf{$R^2$:} Variance in O\_full explained by linear model. Success: $R^2$ > 0.5.
\item \textbf{Scaling factor k:} Slope of regression line O\_full = k $\times O_{10}$. Report mean k and coefficient of variation CV across hypothesis types.
\end{itemize}

\textbf{RQ2 Metrics (Bayesian Error Reduction):}
\begin{itemize}
\item \textbf{Prediction error:} |O\_pred - O\_full| / O\_full (relative error percentage).
\item \textbf{Error reduction:} (Error\_G1 - Error\_G2) / Error\_G1 $\times$ 100\%. Success: >40\%, paired t-test p < 0.05.
\item Sample size: $\geq 10$ hypotheses with Gate 2 data for adequate statistical power.
\end{itemize}

\textbf{RQ3 Metrics (Classification Accuracy):}
\begin{itemize}
\item \textbf{Accuracy:} (TP + TN) / Total. Success: >80\% (24/30 correct), binomial test vs 50\% null, p < 0.05.
\item \textbf{Confusion matrix:} True Positives (non-viable correctly identified), True Negatives (viable correctly identified), False Positives (viable incorrectly stopped), False Negatives (non-viable missed).
\item \textbf{Recall on non-viable:} TP / (TP + FN). Prioritize minimizing false negatives.
\end{itemize}

Statistical significance testing ensures results are not due to chance. For RQ1, Pearson correlation p-value. For RQ2, paired t-test comparing Gate 1 vs Gate 2 errors. For RQ3, binomial test against 50\% null hypothesis.

\subsubsection{4.4 Implementation}

\textbf{Software:} Python 3.8+, scipy.stats for Bayesian updates and statistical tests, numpy for linear regression, matplotlib for visualization.

\textbf{Reproducibility:} All experiments use fixed random seed (42). Synthetic corpus generation is deterministic.

\textbf{Measurement Protocol:} For each hypothesis:
\begin{enumerate}
\item Gate 1: Measure $O_{10}$ on 10 samples (wall-clock time normalized by baseline)
\item Gate 2: Measure $O_{100}$ on 100 samples (subset of hypotheses for RQ2)
\item Gate 3: Measure O\_full on complete synthetic dataset (ground truth)
\end{enumerate}

Synthetic data introduces internal validity risk (perfect linearity artifact). This is addressed by explicitly documenting synthetic nature in results, prioritizing external validation on real corpus as future work, and testing framework mechanics (gate logic, Bayesian updates, statistical tests) rather than claiming generalizability.
